% 04-conformance.tex — Clause 4: Conformance
\chapter{Conformance}
\label{sec:conformance}
\index{conformance}

\section{Conforming implementation}
\label{sec:conformance.impl}
\index{conforming implementation}

\pnum
A \textbf{conforming implementation} shall:

\begin{enumerate}
  \item accept, without diagnostic, every \textit{conforming program};
  \item reject, with at least one diagnostic message, every program that
        violates any constraint stated in this specification;
  \item translate accepted programs such that their observable behavior
        conforms to the semantic rules stated in this specification;
  \item implement all features described as normative in Clauses~5
        through~20;
  \item document every behavior described in this specification as
        \idbehav{the exact value or behavior is implementation-defined}.
\end{enumerate}

\pnum
A conforming implementation may:

\begin{enumerate}
  \item issue additional diagnostic messages (warnings) beyond those
        required by this specification, provided they do not cause a
        conforming program to be rejected;
  \item accept programs through implementation-defined extensions, provided
        those extensions are documented and do not alter the semantics of
        any conforming program;
  \item omit runtime bounds checks when value range analysis
        (\S\,\ref{sec:13-vra}) proves the access is within bounds.
\end{enumerate}

\begin{note}
  The reference implementation of Lain (\textit{lain} compiler, version
  corresponding to this specification) is a conforming implementation.
  Where this specification marks a behavior as implementation-defined, the
  reference implementation's documented behavior takes precedence for
  programs compiled by that implementation.
\end{note}

\section{Conforming program}
\label{sec:conformance.program}
\index{conforming program}

\pnum
A \textbf{conforming program} is a program that:

\begin{enumerate}
  \item consists only of syntax and semantics defined in this
        specification;
  \item does not rely on undefined behavior;
  \item does not rely on unspecified behavior to produce a specific outcome;
  \item does not depend on any implementation-defined extension.
\end{enumerate}

\pnum
A \textbf{strictly conforming program} additionally:

\begin{enumerate}
  \item does not use any behavior documented as implementation-defined,
        even when that behavior is deterministic on a given implementation.
\end{enumerate}

\section{Diagnostic messages}
\label{sec:conformance.diagnostics}
\index{diagnostic message!requirements}

\pnum
A diagnostic message issued for a constraint violation shall include:

\begin{enumerate}
  \item the source file name;
  \item the line number of the offending construct;
  \item the diagnostic code in the form \texttt{[EXXX]} (see Annex~B);
  \item a human-readable description of the violation.
\end{enumerate}

\pnum
A conforming implementation may issue additional diagnostic information
(column number, source context, suggested fix) beyond the minimum required
by \S\,\ref{sec:conformance.diagnostics}.

\begin{note}
  The reference implementation additionally prints the source line and a
  caret (\texttt{\^{}}) pointing to the column of the offending token.
\end{note}

\section{Implementation limits}
\label{sec:conformance.limits}
\index{implementation limits}

\pnum
A conforming implementation shall support at least the following minimum
translation limits.  Programs that require more than these limits invoke
\idbehav{the behavior when a translation limit is exceeded}.

\begin{longtable}{@{}lll@{}}
\toprule
\textbf{Resource} & \textbf{Minimum} & \textbf{Clause} \\
\midrule
\endfirsthead
\multicolumn{3}{l}{\textit{(continued)}}\\
\toprule
\textbf{Resource} & \textbf{Minimum} & \textbf{Clause} \\
\midrule
\endhead
\bottomrule
\endlastfoot
Identifier length (bytes)                      & 255         & \S\,\ref{sec:05-lexical} \\
Block nesting depth                            & 64          & \S\,\ref{sec:09-statements} \\
Function parameters                            & 127         & \S\,\ref{sec:10-declarations} \\
Array elements                                 & $2^{31}-1$  & \S\,\ref{sec:07-types} \\
Module import depth                            & 32          & \S\,\ref{sec:16-modules} \\
\lain{case} arms per match expression          & 256         & \S\,\ref{sec:15-match} \\
Struct fields                                  & 256         & \S\,\ref{sec:07-types} \\
Enum variants                                  & 256         & \S\,\ref{sec:07-types} \\
Comptime evaluation depth                      & 256         & \S\,\ref{sec:14-generics} \\
Scope nesting depth                            & 64          & \S\,\ref{sec:06-basics} \\
\end{longtable}

\section{Safe and unsafe fragments}
\label{sec:conformance.safe-unsafe}
\index{safe fragment}
\index{unsafe fragment}

\pnum
The \textbf{safe fragment} of a Lain program consists of all code that does
not appear inside an \lain{unsafe} block and does not use raw pointer
operations.

\pnum
A conforming implementation shall guarantee, for any well-formed program
in the safe fragment, that execution produces no:

\begin{itemize}
  \item use of a variable after its ownership has been transferred (use-after-move);
  \item double transfer of ownership of the same value;
  \item access to an uninitialized variable;
  \item violation of the read-write lock invariant
        (\S\,\ref{sec:terms.ownership});
  \item buffer overrun that has been statically proven safe by value range
        analysis.
\end{itemize}

\pnum
The \textbf{unsafe fragment} comprises all code inside \lain{unsafe}
blocks.  Within the unsafe fragment, the programmer is responsible for
upholding the memory safety invariants listed in
\S\,\ref{sec:18-unsafe}; the implementation shall not be required to
verify them.

\begin{constraint}
  The constraints of the ownership and linearity system
  (\S\,\ref{sec:11-ownership}) are \emph{not} suspended inside
  \lain{unsafe} blocks.  Only raw pointer operations and the bounds-check
  requirement are relaxed.  An implementation shall continue to enforce
  \diagcode{E001}--\diagcode{E005} inside \lain{unsafe} blocks.
\end{constraint}
